\documentclass{article}
\usepackage[T2A]{fontenc} % russian letters
\usepackage{graphicx} % Required for inserting images
\usepackage{tikz}
\usepackage{geometry}
\usepackage{multirow}
\usepackage{colortbl}

\geometry{
    top=1cm,
    bottom=2cm,
    left=4cm,
    right=4cm
}

% Определение стиля для узлов в кружочке
\tikzset{
  nodeStyle/.style={circle, draw, minimum size=10mm, inner sep=0pt},
}

\renewcommand{\tablename}{Таблица}

\title{Лабароторная работа №2}
\author{Лебедев Денис и Кветко Матвей, 2 курс 9 группа}
\date{26 сентября 2023 г.}

\begin{document}

\maketitle
\section{Теоретические сведения}
\hspace*{1cm}Доспустим у нас есть какой-либо граф, тогда из него можно выделите следующие подграфы:
\begin{itemize}
    \item \textbf{Остовное дерево} - подграф, который включает в себя все вершины исходного графа и является деревом;
    \item \textbf{Покрывающее дерево} - подграф ориентированного или неориентированного графа, который содержит все вершины графа и использует наименьшее количество рёбер;
    \item \textbf{Корневое дерево} - это дерево, в котором одна из вершин называется корнем, и каждое ребро направлено от корня к листьям.
\end{itemize}
\hspace*{1cm}Для каждой вершины корневого дерева мы можем найти глубину (расстояние от корня до этой вершины), список предков и направление, что поможет при создании таблицы представления корневого дерева в памяти компьютера. Также для этого необходимо знать, как строить \textbf{династический обход} - метод обхода дерева, который позволяет посещать вершины дерева в определённом порядке, учитывая их иерархическую структуру.

\section{Изначальный граф}
\hspace*{1cm}На рисунке 1 представлен граф \textit{S = (I,U), I} = \{1, 2, 3, 4, 5, 6, 7\}, \textit{U} = \{\{1, 5\}, \{1, 6\}, \{2, 1\}, \{2, 3\}, \{2, 6\}, \{2, 7\}, \{3, 1\}, \{4, 3\}, \{5, 2\}, \{5, 4\}, \{6, 4\}, \{6, 3\}\};
\begin{center}
\begin{tikzpicture}
  \node[nodeStyle] (1) at (1,4) {1};
  \node[nodeStyle] (2) at (0,2) {2};
  \node[nodeStyle] (3) at (2,0) {3};
  \node[nodeStyle] (4) at (4,0) {4};
  \node[nodeStyle] (5) at (6,2) {5};
  \node[nodeStyle] (6) at (5,4) {6};
  \node[nodeStyle] (7) at (3,5) {7};

  \draw[->, >=stealth, ultra thick] (1) -- (5);
  \draw[->, >=stealth, ultra thick] (1) -- (6);
  \draw[->, >=stealth, ultra thick] (2) -- (1);
  \draw[->, >=stealth, ultra thick] (2) -- (3);
  \draw[->, >=stealth, ultra thick] (2) -- (6);
  \draw[->, >=stealth, ultra thick] (2) -- (7);
  \draw[->, >=stealth, ultra thick] (3) -- (1);
  \draw[->, >=stealth, ultra thick] (4) -- (3);
  \draw[->, >=stealth, ultra thick] (5) -- (2);
  \draw[->, >=stealth, ultra thick] (5) -- (4);
  \draw[->, >=stealth, ultra thick] (6) -- (4);
  \draw[->, >=stealth, ultra thick] (6) -- (3);

  \node[below] at (3, -1) {Рисунок 1. Граф \textit{S = (I,U)}};
\end{tikzpicture}
\end{center}

\newgeometry{top=7cm, bottom=2cm}
\newpage
\section{Лебедев Денис}
\hspace*{1cm} На рисунке 2 представлено корневое дерево данного графа с корнем 3 и соответствующее остовное дерево. В таблице 1 отражено представление данного корневого дерева в памяти компьютера.\\\\
\begin{tikzpicture}
  \node[nodeStyle] (1) at (6,2) {1};
  \node[nodeStyle] (2) at (2,2) {2};
  \node[nodeStyle] (3) at (3,4) [fill=gray!50] {3};
  \node[nodeStyle] (4) at (4,2) {4};
  \node[nodeStyle] (5) at (1,0) {5};
  \node[nodeStyle] (6) at (0,2) {6};
  \node[nodeStyle] (7) at (3,0) {7};

  \draw[->, >=stealth, ultra thick] (2) -- (5);
  \draw[->, >=stealth, ultra thick] (2) -- (7);
  \draw[->, >=stealth, ultra thick] (3) -- (1);
  \draw[->, >=stealth, ultra thick] (3) -- (2);
  \draw[->, >=stealth, ultra thick] (3) -- (4);
  \draw[->, >=stealth, ultra thick] (3) -- (6);

  \node[nodeStyle] (1) at (14,2) {1};
  \node[nodeStyle] (2) at (10,2) {2};
  \node[nodeStyle] (3) at (11,4) {3};
  \node[nodeStyle] (4) at (12,2) {4};
  \node[nodeStyle] (5) at (9,0) {5};
  \node[nodeStyle] (6) at (8,2) {6};
  \node[nodeStyle] (7) at (11,0) {7};

  \draw[-, >=stealth, ultra thick] (2) -- (5);
  \draw[-, >=stealth, ultra thick] (2) -- (7);
  \draw[-, >=stealth, ultra thick] (3) -- (1);
  \draw[-, >=stealth, ultra thick] (3) -- (2);
  \draw[-, >=stealth, ultra thick] (3) -- (4);
  \draw[-, >=stealth, ultra thick] (3) -- (6);

  \node[below] at (7, -1) {Рисунок 2. Корневое дерево с корнем 3 (слева), соответствующее остовное дерево (справа)};
\end{tikzpicture}

\begin{table} [h]
  \centering
  \caption{Представление корневого дерева в памяти компьютера}
  \begin{tabular}{|c|c|c|c|c|c|c|c|}
    \hline
    \cellcolor{gray!75}Структуры представления & \multicolumn{7}{|c|}{\cellcolor{gray!50} Номер узла} \\
    \arrayrulecolor{gray!75}\cline{1-2}
    \arrayrulecolor{black} \cline{2-8}
    \cellcolor{gray!75}корневого дерева & \cellcolor{gray!35}1 & \cellcolor{gray!35}2 & \cellcolor{gray!35}3 & \cellcolor{gray!35}4 & \cellcolor{gray!35}5 & \cellcolor{gray!35}6 & \cellcolor{gray!35}7 \\
    \hline
    \cellcolor{gray!50}\textit{pred[i]} & 3 & 3 & 0 & 3 & 2 & 3 & 2 \\
    \hline
    \cellcolor{gray!50}\textit{depth[i]} & 1 & 1 & 0 & 1 & 2 & 1 & 2 \\
    \hline
    \cellcolor{gray!50}\textit{dir[i]} & 1 & -1 & 0 & -1 & -1 & -1 & 1 \\
    \hline
    \cellcolor{gray!50}\textit{d[i]} & 4 & 7 & 1 & 2 & 6 & 3 & 5 \\
    \hline
  \end{tabular}
\end{table}

\newpage

\section{Кветко Матвей}
\hspace*{1cm} На рисунке 3 представлено корневое дерево данного графа с корнем 7, а также соответствующее остовное дерево. В таблице 2 приведено представление этого корневого дерева в памяти компьютера.\\\\
\begin{tikzpicture}
  \node[nodeStyle] (1) at (2,4) {1};
  \node[nodeStyle] (2) at (5,5) {2};
  \node[nodeStyle] (3) at (5,3) {3};
  \node[nodeStyle] (4) at (3,1) {4};
  \node[nodeStyle] (5) at (8,4) {5};
  \node[nodeStyle] (6) at (1,2.5) {6};
  \node[nodeStyle] (7) at (5,7) [fill=gray!50] {7};

  \draw[->, >=stealth, ultra thick] (1) -- (6);
  \draw[->, >=stealth, ultra thick] (2) -- (1);
  \draw[->, >=stealth, ultra thick] (2) -- (3);
  \draw[->, >=stealth, ultra thick] (2) -- (5);
  \draw[->, >=stealth, ultra thick] (6) -- (4);
  \draw[->, >=stealth, ultra thick] (7) -- (2);

  \node[nodeStyle] (1) at (10,4) {1};
  \node[nodeStyle] (2) at (13,5) {2};
  \node[nodeStyle] (3) at (13,3) {3};
  \node[nodeStyle] (4) at (11,1) {4};
  \node[nodeStyle] (5) at (16,4) {5};
  \node[nodeStyle] (6) at (9,2.5) {6};
  \node[nodeStyle] (7) at (13,7) {7};

  \draw[-, >=stealth, ultra thick] (1) -- (6);
  \draw[-, >=stealth, ultra thick] (2) -- (1);
  \draw[-, >=stealth, ultra thick] (2) -- (3);
  \draw[-, >=stealth, ultra thick] (2) -- (5);
  \draw[-, >=stealth, ultra thick] (6) -- (4);
  \draw[-, >=stealth, ultra thick] (7) -- (2);

  \node[below] at (8, 0) {Рисунок 3. Корневое дерево с корнем 7 (слева), соответстующее остовное дерево (справа)};
\end{tikzpicture}

\begin{table} [h]
  \centering
  \caption{Представление корневого дерева в памяти компьютера}
  \begin{tabular}{|c|c|c|c|c|c|c|c|}
    \hline
    \cellcolor{gray!75}Структуры представления & \multicolumn{7}{|c|}{\cellcolor{gray!50} Номер узла} \\
    \arrayrulecolor{gray!75}\cline{1-2}
    \arrayrulecolor{black} \cline{2-8}
    \cellcolor{gray!75}корневого дерева & \cellcolor{gray!35}1 & \cellcolor{gray!35}2 & \cellcolor{gray!35}3 & \cellcolor{gray!35}4 & \cellcolor{gray!35}5 & \cellcolor{gray!35}6 & \cellcolor{gray!35}7 \\
    \hline
    \cellcolor{gray!50}\textit{pred[i]} & 2 & 7 & 2 & 6 & 2 & 1 & 0 \\
    \hline
    \cellcolor{gray!50}\textit{depth[i]} & 2 & 1 & 2 & 4 & 2 & 3 & 0 \\
    \hline
    \cellcolor{gray!50}\textit{dir[i]} & 1 & -1 & 1 & 1 & -1 & 1 & 0 \\
    \hline
    \cellcolor{gray!50}\textit{d[i]} & 6 & 5 & 1 & 7 & 3 & 4 & 2 \\
    \hline
  \end{tabular}
\end{table}

\restoregeometry

\end{document}
